\documentclass[11pt,a4paper]{letter}
\usepackage[T1]{fontenc}
\usepackage[utf8]{inputenc}
\usepackage[left=2.5cm,right=2.5cm,top=2.5cm,bottom=2.5cm]{geometry}
\usepackage{hyperref}

\signature{Rodrigo Almeida de Oliveira, M.Sc. \\ PhD candidate, PPGIa \\ Pontif\'icia Universidade Cat\'olica do Paran\'a (PUCPR)}
\address{Rodrigo Almeida de Oliveira, M.Sc. \\ Pontif\'icia Universidade Cat\'olica do Paran\'a \\ Programa de P\'os-Gradua\c c\~ao em Inform\'atica (PPGIa) \\ Curitiba, PR, Brazil \\ rodrigo1.almeida@pucpr.edu.br}

\begin{document}

\begin{letter}{Editor-in-Chief \\ \emph{Information and Software Technology} \\ Elsevier}

\opening{Dear Editors,}

On behalf of all co-authors (Rodrigo Almeida de Oliveira, Juliano de Paulo
Ribeiro, and Prof.\ Edson Emilio Scalabrin), I am pleased to submit our
manuscript titled
\textbf{``Process Mining and Stochastic Modeling for Software Process
Forecasting: A Systematic Literature Review with an LLM-Assisted Protocol''}
for consideration in \emph{Information and Software Technology}.

\medskip\noindent\textbf{Novelty and contribution.}
This manuscript reports a systematic literature review (SLR) at the
intersection of process mining (PM), stochastic modeling, and machine
learning applied to software process forecasting. The contribution is
threefold:
\begin{enumerate}
\item \textbf{Methodological:} we operationalise an LLM-assisted screening
      protocol with cross-model verification (Claude Haiku 4.5 as primary
      screener, Claude Sonnet 4.6 as verifier), achieving substantial
      Cohen's $\kappa$ (binary $\kappa = 0.695$ for title/abstract,
      $0.694$ for full-text) on a stratified random 20\% sample. The
      protocol extends the recent work of Khraisha et~al.~(2024) and
      Syriani et~al.~(2023) to the software-process-mining domain at a
      scale of 6{,}147 deduplicated papers across two corpus tiers.
\item \textbf{Empirical:} we confirm 404 primary studies (169 working-set
      + 212 + 23 auxiliary) covering 1995--2026, scored against an
      8-criterion quality rubric (Dyba and Dingsoyr, 2008). The combined
      mean QA score is $5.00 \pm 1.42$ (315 papers, 82.7\%, retained for
      synthesis).
\item \textbf{Synthetic:} we report five empirically grounded findings
      (F1--F5), propose the SPMF taxonomy (Software Process Mining and
      Forecasting) along three orthogonal dimensions, and articulate a
      four-direction research agenda (RA1--RA4). The IC-combination
      analysis shows that only 1 of 404 papers (0.3\%) jointly satisfies
      process mining, stochastic modeling, and forecasting --- evidence
      that the gap addressed by our PATHCAST proposal is structural
      rather than an artefact of the search queries.
\end{enumerate}

\medskip\noindent\textbf{Fit with IST scope.}
The manuscript squarely fits IST's scope on empirical software engineering
methods, mining software repositories, and process modeling. It complements
the journal's growing body of evidence-based SLRs while contributing to
the methodological literature on LLM-assisted reviews.

\medskip\noindent\textbf{Reproducibility.}
A complete replication package is deposited at Zenodo
(\href{https://doi.org/10.5281/zenodo.15719919}{DOI~10.5281/zenodo.15719919}),
including all search strings, screening prompts, JSON schemas, decision
CSVs at every pipeline stage, the cross-model $\kappa$ scripts, and the
structured extraction tables.

\medskip\noindent\textbf{Funding.}
This study was financed in part by CAPES (Finance Code 001) through a
doctoral scholarship granted to the corresponding author within the
Graduate Program in Informatics (PPGIa) at PUCPR.

\medskip\noindent\textbf{Prior submission and conflicts.}
The manuscript is original work, has not been published elsewhere, and is
not under consideration by any other journal. All authors declare no
competing financial interests. The use of generative AI tools in the
research is declared in the manuscript per Elsevier policy. AI tools are
\emph{not} listed as authors and bear no responsibility for editorial
judgements.

\medskip\noindent\textbf{Suggested reviewers.}
\begin{itemize}
\item \textbf{Wil M.\ P.\ van der Aalst} --- TU Eindhoven, foundational PM
      researcher (\texttt{w.m.p.v.d.aalst@tue.nl})
\item \textbf{Marcelo Sergio Pavlovic} --- LASIGE, FCUL, PM in software
      engineering
\item \textbf{Maria Teresa Baldassarre} --- University of Bari, empirical
      SE and SLRs
\item \textbf{Eugene Syriani} --- University of Montreal, LLM-assisted SLR
      methodology
\item \textbf{Tore Dyb{\aa}} --- SINTEF / University of Oslo, evidence-based
      SE
\end{itemize}

\medskip\noindent\textbf{Reviewers to avoid.}
None on basis of co-authorship within the past five years. Specific
exclusions: thesis advisors and committee members at PUCPR.

\closing{Sincerely,}

\end{letter}

\end{document}
